% -----------------------------------------------------------------------------
% This file is automatically generated.
% Do not edit manually.
% Source: notebooks/support/regression-analysis/support.Rmd
% -----------------------------------------------------------------------------


\noindent A common rule of thumb is that VIF values greater
than 10 indicate substantial multicollinearity.
Since none of the VIF values exceeds this
threshold, we conclude that there is no evidence
of serious variance inflation in \ttblue{mod.1}.
For models that include interaction terms, such as
\ttblue{mod.3}, we use the \ttblue{vif()} function with the
additional option \ttblue{type = "predictor"}. This reports
the Generalized Variance Inflation Factor (GVIF) and its
adjusted value, \(GVIF^{1/(2df)}\).

\par\addvspace{\topsep}
\begingroup
\fvset{listparameters={%
  \setlength{\topsep}{0pt}%
  \setlength{\partopsep}{0pt}%
  \setlength{\parsep}{0pt}%
  \setlength{\itemsep}{0pt}%
}}
\begin{Verbatim}[breaklines=true,formatcom=\color{ada_blue}]
model_vif('ussteamco_mod_3', USSteamCoEstim2)
\end{Verbatim}
\begin{Verbatim}[breaklines=true,formatcom=\color{ada_light_blue}]
               GVIF Df GVIF^(1/(2*Df))
production 203.6856  3        2.425642
heatDD     192.5840  3        2.403090
coolDD     226.3290  3        2.468634
summer       1.0000  7        1.000000
                       Interacts With   Other Predictors
production                     summer     heatDD, coolDD
heatDD                         summer production, coolDD
coolDD                         summer production, heatDD
summer     production, heatDD, coolDD               --
\end{Verbatim}
\endgroup
\par\addvspace{\topsep}
\noindent The adjusted GVIF places predictors with
different degrees of freedom on a comparable
scale. For a predictor with one degree of
freedom, the adjusted GVIF equals the square root
of the ordinary VIF. Consequently, the
conventional VIF threshold of 10 corresponds to
an adjusted GVIF of approximately \(\sqrt{10}
\approx 3.16\). Since none of the reported values
exceeds 3.16, we conclude that \ttblue{mod.3}
does not exhibit problematic multicollinearity.
